\section{Giải pháp kết nối Laptop, Desktop và Smartphone}
\label{sec:device-connectivity}

Trong mô hình Infrastructure, các thiết bị đầu cuối kết nối tới cùng điểm truy
cập (access point -- AP) bằng định danh mạng (SSID) và phương thức xác thực
được cấu hình thống nhất. Sau khi liên kết lớp vô tuyến, thiết bị nhận cấu hình
địa chỉ IP tự động qua DHCP hoặc sử dụng cấu hình tĩnh phù hợp với mạng.

\subsection{Máy tính xách tay (Laptop)}
Kiểm tra bộ điều hợp Wi-Fi đã bật và hỗ trợ băng tần, chuẩn bảo mật của AP. Chọn
SSID, nhập thông tin xác thực, xác nhận địa chỉ IP, cổng mặc định và máy chủ DNS;
kiểm tra kết nối mạng nội bộ và Internet.

\subsection{Máy tính để bàn (Desktop)}
Máy tính để bàn có thể dùng card Wi-Fi PCIe/USB hoặc kết nối Ethernet tới LAN.
Nếu dùng Wi-Fi, thực hiện ghép nối như Laptop; nếu dùng Ethernet, bảo đảm cổng
LAN của router/AP hoạt động và cấu hình IP tương thích với cùng mạng con.

\subsection{Điện thoại thông minh (Smartphone)}
Trên Android hoặc iOS, bật Wi-Fi, chọn SSID và xác thực theo chính sách của
mạng. Với mạng doanh nghiệp, nhập thông tin 802.1X theo hướng dẫn quản trị viên.
Không chia sẻ mật khẩu qua kênh không an toàn; ưu tiên WPA2-AES hoặc WPA3 khi
thiết bị hỗ trợ.

\subsection{Kết nối ngang hàng Ad-hoc}
Trong mạng Ad-hoc (IBSS), các trạm tương thích kết nối trực tiếp, không qua AP.
Khả năng hỗ trợ tùy hệ điều hành, trình điều khiển và phần cứng; Wi-Fi Direct
là một cơ chế ngang hàng phổ biến trên thiết bị hiện đại nhưng không đồng nhất
hoàn toàn với IBSS.
